\section{Prototipagem}



\subsection{Energy Dispersal}
\label{sec:energy}
O fluxo de dados de link de retorno MF-TDMA (Multi-Frequency Time-Division multiple-access) é organizado em pacotes de dados . Para cumprir os Regulamentos de Radiocomunicações ITU e para garantir transições binárias adequadas o fluxo de bits em uma sequencia deve ser misturado para reduzir o probabilidade de sequências de bits monotônicas. Uma sequência de pseudo aleatória de números binários (PRBS - Pseudo-Random Binary Sequence)  deve ser usada conforme especificado pela expressão polinomial [\ref{eq:polyEnergy}], tal expressão surge de descrição matemática do diagrama de blocos da figura \ref{fig:energy} e também da quantidade de bits utilizada \cite{Hideo}.

\begin{equation}\label{eq:polyEnergy}
1+x^{14}+x^{15}
\end{equation}


Os dados são randomizados usando o registrador Linear Feedback Shift Register (LFSR) de 15 registros mostrado na figura \ref{fig:energy} para distribuir a distribuição de uns e zeros. O gerador realiza a adição do módulo-2 dos dados com a pseudo-aleatória sequência. O conteúdo inicial dos registros SR-1 a SR-15 é dado na tabela \ref{tab:InitRand}. O primeiro bit da sequência pseudo-aleatória  é será adicionado módulo 2 com o primeiro bit do fluxo de bits de dados em série, isto é, o primeiro bit após o preâmbulo de explosão.
O gerador é redefinido para o conteúdo inicial antes de processar uma carga útil do pacote de dados.
\begin{table}[!ht]
      \begin{center}
         \caption{Valor Inicial do Randomizer}
         \label{tab:InitRand}
        \includegraphics[width=3.5in]{Images/initRandomizer.png}
      \end{center}
\end{table}

 \begin{figure}[!ht]
        \begin{center}
            \includegraphics[width=3.5in]{Images/LFSR.png}
            \caption{{Energy Dispersal}}
            \label{fig:energy}
        \end{center}
 \end{figure}

O Frame PDU, um pacote de informação que une dados de controle de protocolo e dados do usuário, deve ser emitido para o energy dispersal como uma sequência de bits com conteúdo que segue:
\begin{enumerate}{}
	\item Os campos fixos das estruturas de dados devem ser emitidos na ordem em que estão listados, de cima para baixo.
	\item Uma sequência continua de campos como uma iteração de um tipo de campo deve ser emitida na ordem da iteração.
	\item Campos relocáveis de uma estrutura de dados flexível podem ser emitidos em qualquer ordem, se não explicitamente indicado de outra forma.
	\item Os bits reservados associados a um campo devem ser emitidos antes dos bits de informação do mesmo campo.
	\item Os campos multi-bit serão emitidos em ordem de bits começando com o bit considerado mais significativo e terminando com o pouco considerado menos significativo.
\end{enumerate}

\subsection{Playload CRC16}
\label{sec:crc16}
O CRC (cyclic redundancy check) é um código de detecção de erros usado comumente em sistemas digitais e em dispositivos de armazenamento com intuito de detectar mudanças errôneas no dados brutos. Blocos de dados que entram nesses sistemas recebem um valor de verificação curto, com base no restante de uma divisão polinomial de seus conteúdos. Na recuperação, o cálculo é repetido e, caso os valores de verificação não coincidam, uma ação corretiva pode ser tomada contra a corrupção de dados \cite{Bill}.  

O modulo CRC16 deve ser calculado para a saída de sequência de bits do modulo de Energy Dispersal. Também deve ser calculado sobre o Frame PDU (exclusivo do próprio CRC). O polinômio utilizado é um polinômio padrão e pode ser visto na \ref{eq:2}

\begin{equation}\label{eq:2}
x^{16}+x^{15}+x^2+1
\end{equation}

O CRC é o restante da divisão da carga útil de burst pelo polinômio.
O CRC16 deve ser equivalente ao calculado por um circuito, como mostrado na figura \ref{sec:crc16}. As células do registrador de deslocamento devem ser inicializado para 0 antes do início do cálculo. Primeiro, os interruptores estão na posição "A" e a palavra de dados é deslocada e transmitido simultaneamente. 

Após o último bit PDU do quadro, os interruptores são movidos para a posição "B", e o o conteúdo do registrador de deslocamento é transmitido, começando com o bit no final do registrador. Este é o diagrama de blocos do CRC16 figura \ref{fig:crc16}.

\begin{figure}[!ht]
      \begin{center}
        \includegraphics[width=3.5in]{Images/CRC16.png}
        \caption{{Diagrama de blocos do CRC16 }}
        \label{fig:crc16}
      \end{center}
\end{figure}
